\begin{pastexam}[label={px:ccu-111-swa-04}]{中正資工所甲組 111 學年度 第 12 題}
（單選）

(3\%) (\quad) Given an unsorted array containing the following numbers. You should construct a
binary tree from the array and convert it into a max-heap.
\begin{center}45, 83, 7, 61, 12, 99, 44, 77, 14, 29\end{center}
What does the array look like after the second phase (continuously output the first two largest
numbers) of performing a heap sorting?

a.~77, 61, 45, 14, 29, 7, 44, 12 \qquad b.~45, 7, 61, 12, 44, 77, 14, 29\\
c.~77, 61, 44, 45, 29, 7, 12, 14 \qquad d.~77, 61, 44, 45, 12, 7, 29, 14\\
e.~None of the above
\end{pastexam}

\begin{pastsoln}
答案是 a。

題目說「把陣列看成一棵二元樹，再轉成最大堆積」，這是資結書第~8.2~節由下往上的建堆積（虛擬碼~8.3）：
陣列不動，從最後一個內部節點往回，對每個節點做向下調整。 堆排序（資結書虛擬碼~14.13）第一階段建
堆積，第二階段每次把根（最大值）和堆積最後一格交換、堆積縮小一格、再從根往下調整。 題目要的是輸出
兩個最大值之後，堆積裡剩下的 $8$ 格。 下面陣列從第 $1$ 格算，第 $i$ 格的子節點是第 $2i$、$2i + 1$ 格。

\textbf{第一階段：建堆積。} $n = 10$，從第 $\lfloor 10/2 \rfloor = 5$ 格往回做到第 $1$ 格。 第 $1$ 格的 45 和
99 交換後到了第 $3$ 格，子節點 7、44 都比 45 小，停下。
\begin{center}
\small
\begin{tabular}{cll}
\toprule
調整第 $i$ 格 & 動作 & 陣列 \\
\midrule
（開始） & & 45, 83, 7, 61, 12, 99, 44, 77, 14, 29 \\
$i = 5$（12） & 和子節點 29 交換 & 45, 83, 7, 61, \textbf{29}, 99, 44, 77, 14, \textbf{12} \\
$i = 4$（61） & 和較大的子節點 77 交換 & 45, 83, 7, \textbf{77}, 29, 99, 44, \textbf{61}, 14, 12 \\
$i = 3$（7） & 和較大的子節點 99 交換 & 45, 83, \textbf{99}, 77, 29, \textbf{7}, 44, 61, 14, 12 \\
$i = 2$（83） & 子節點 77、29 較小，不動 & 45, 83, 99, 77, 29, 7, 44, 61, 14, 12 \\
$i = 1$（45） & 和 99 交換後停 & \textbf{99}, 83, \textbf{45}, 77, 29, 7, 44, 61, 14, 12 \\
\bottomrule
\end{tabular}
\end{center}

\textbf{第二階段：輸出兩個最大值。}
\begin{enumerate}
\item 99 和最後一格的 12 交換，99 離開堆積：12, 83, 45, 77, 29, 7, 44, 61, 14 $\mid$ 99。 12 往下：和
83 交換（第 2 格），再和 77 交換（第 4 格），再和 61 交換（第 8 格），到底。 得
83, 77, 45, 61, 29, 7, 44, 12, 14 $\mid$ 99。
\item 83 和第 9 格的 14 交換：14, 77, 45, 61, 29, 7, 44, 12 $\mid$ 83, 99。 14 往下：和 77 交換（第 2 格），
再和 61 交換（第 4 格）；第 4 格的子節點只有第 8 格的 12，比 14 小，停。 得
\[
77,\ 61,\ 45,\ 14,\ 29,\ 7,\ 44,\ 12 \ \mid\ 83,\ 99.
\]
\end{enumerate}
堆積裡的 $8$ 格就是 a。

逐項檢查：
\begin{itemize}
\item a：上面的結果。 對。
\item b：45, 7, 61, \ldots{} 的根 45 比子節點 61 小，不是最大堆積。 錯。
\item c：這是把十個數\textbf{一個一個插入}空堆積（每次往上調整）建出來的堆積
99, 77, 83, 61, 29, 7, 44, 45, 14, 12，再輸出兩次的結果。 題目要的是把整個陣列當成樹再調整，兩種建法
得到的堆積不同。 錯。
\item d：八個數對，也是合法的最大堆積（每一格都不小於它的子節點），但不是這個程序算出來的；
和 a 比，第 3 格就不同（a 是 45，d 是 44）。 錯。
\item e：a 是對的。 錯。
\end{itemize}
\end{pastsoln}
